\BeleskeID{02-s015}
% element: 02-s015:e01
% element: 02-s015:e02
% element: 02-s015:d01

Kada uslov izgovaramo koristeći veznik I, najpre prepoznajemo kombinaciju ulaza: prava vrednost promenljive iznosi 0, pa njena komplementna vrednost iznosi 1. Zbog toga sva tri literala u proizvodu \(\bar C\bar B\bar A\) imaju vrednost 1 za red 000. Taj proizvod opisuje ispunjenost uslova, ali nije sam činilac izlazne funkcije koju formiramo iz redova sa nulom.

Negacija proizvoda daje traženi činilac sa aktivnom nulom. U bilo kom drugom redu bar jedna prava ulazna promenljiva ima vrednost 1, pa je zbir \(C+B+A\) jednak 1. Tako činilac izdvaja samo jedan red nulom. Dualni NI i ILI opis odnose se na isti logički uslov; pri čitanju moramo pratiti i aktivni nivo, a ne samo veznik u rečenici.
